\documentclass{article}
\usepackage{readmasters}
\begin{document}

\origpage{298}

\section*{Zur Theorie der algebraischen Functionen mehrerer complexer Variabeln.}

\subsection*{Von Max Nöther.}

In der Theorie der algebraischen Functionen \emph{einer} complexen Variabeln hat Riemann eine Klassification der Gleichungen aufgestellt${}^{1)}$, welche diese Functionen definiren. Das Charakteristische einer Klasse ist, dass sich alle derselben angehörigen Gleichungen eindeutig durch rationale Substitutionen in einander transformiren lassen. Wenn $f = o$ die homogene Gleichung einer Curve $n^{\text{ter}}$ Ordnung mit $d$ Doppel- und $r$ Rückkehrpunkten${}^{2)}$, so ist $p = \frac{n-1.\; n-2}{1.\; 2.} - d - r$ die Klassenzahl

\textbf{1)} Crelle's Journal, Bd.\ 54.

\textbf{2)} Clebsch und Gordan, Theorie der Abelschen Funct.

\origpage{299}
der Klasse, zu welcher $f = o$ gehört, und die Gleichheit dieser Zahl für zwei Curven ist die nothwendige und zugleich hinreichende Bedingung für das eindeutige Entsprechen der beiden Curven. Erst vor kurzer Zeit hat Herr Clebsch dieses Theorem auf Functionen zweier Variabeln erweitert${}^{1)}$. Wenn sich zwei Flächen, die nur die gewöhnlichen Singularitäten, Doppel- und Rückkehrcurve, besitzen, rational in einander transformiren lassen, so haben dieselben hiernach eine Klassenzahl $p$ gemeinschaftlich, die Herr Clebsch als das \emph{Geschlecht} der Fläche bezeichnet. Das Geschlecht $p$ einer Fläche $n^{\text{ter}}$ Ordnung $f = o$ ist die Anzahl der in einer Fläche $(n-4)^{\text{ter}}$ Ordnung noch unbestimmt bleibenden Coefficienten, wenn man diese Fläche durch die Doppel- und Rückkehrcurve von $f = o$ hindurchgehen lässt.

Ich beabsichtige in dieser Mittheilung, diese Theoreme nach mehreren Seiten hin auszudehnen. Einmal werde ich auch die Flächen mit höheren konischen Knotenpunkten und höheren vielfachen Curven in den Kreis der Betrachtung hereinziehen; zweitens werde ich sogleich diese Klasseneintheilung für die algebraischen Gebilde von beliebig vielen Dimensionen geben, in dem nämlichen Umfange, wie dies eben für Flächen angedeutet wurde, und drittens werde ich nachweisen, dass ausser der Gleichheit der als „Geschlecht“ eines Gebildes bezeichneten Zahl $p$ noch weitere Bedingungen für das eindeutige Entsprechen zweier Gebilde existiren, dass man nämlich innerhalb jeder durch ein $p$ charakterisirten Klasse von Gebilden von $2r$ Dimensionen noch eine $(r-1)$ fach unendliche Reihe von

\textbf{1)} Comptes Rendus vom 21.\ Dec.\ 1868, p.\ 1238.

\origpage{300}
Klassen statuiren muss, deren charakteristische Zahlen ganz ähnliche Bedeutung haben, wie $p$ selbst.

\section*{I.}

Eine algebraische Function $s$ von $r$ unabhängig veränderlichen complexen Grössen $x_1, x_2, \ldots x_r$ wird definirt durch eine algebraische Gleichung $n^{\text{ten}}$ Grades
\[ f(s, x_1, x_2, \ldots x_r) = o, \]
die eine $2r$ fach unendliche Mannigfaltigkeit vorstellt. Zur Abkürzung werde ich im Folgenden eine $2h$ fach unendliche Mannigfaltigkeit als eine $\infty^{2h}$ bezeichnen. $f = o$ kann vielfach zu zählende $\infty^{2h}$ enthalten, für $h = o$ bis $h = r - 1$, wie eine Fläche vielfache Curven und Knotenpunkte enthalten kann. Wir setzen jedoch voraus, dass diese vielfachen auf $f = o$ liegenden Mannigfaltigkeiten nicht selbst wieder specielle Singularitäten enthalten. die denen analog wären, welche bei einer Fläche durch das Zusammenfallen von Tangentenebenen längs einer vielfachen Curve oder durch den Uebergang eines konischen Knotenpunktes in einen planaren eintreten. Weiter setzen wir immer voraus, dass die Gleichung $f = o$ irreducibel sei.

\emph{Theorem 1:} „Das \emph{Geschlecht} $p$ dieser Mannigfaltigkeit $f = o$ ist gleich der Anzahl der in einer $\infty^{2r}$ vom Grade $n - r - 2$, $\theta = o$, noch unbestimmt bleibenden Constanten, wenn $\theta = o$ gezwungen wird, durch die vielfachen auf $f = o$ liegenden Mannigfaltigkeiten hindurchzugehen, und zwar durch jede $\mu$fach zählende
\[
\begin{array}{lllll}
\infty^{2(r-1)} & \text{auf} & f = o & (\mu - 1) & \text{mal,}\\
\infty^{2(r-2)} & \text{„} & \text{„} & (\mu - 2) & \text{„}\quad\text{etc.}
\end{array}
\]
endlich durch eine $\mu$fache Curve von $f$ $(\mu - r$

\origpage{301}
$+ 1)$ mal, und einen $\mu$fachen konischen Knotenpunkt $(\mu - r)$ mal.

Diese Zahl $p$ hat die Eigenschaft, für jede andere $\infty^{2r}$, $\varphi = o$, welche $f = o$ im Allgemeinen Punkt für Punkt eindeutig entspricht, die nämliche zu sein, so dass $p$, das Geschlecht von $f = o$, die Klassenzahl ist für die Klasse, zu der $f = o$ gehört. Die Gleichheit der Zahl $p$ ist das nothwendige Kriterium für die Möglichkeit einer eindeutigen rationalen Transformation zweier $\infty^{2r}$ in einander.“

\section*{II.}

Bei dem eindeutigen Entsprechen höherer Mannigfaltigkeiten kommen Eigenthümlichkeiten vor, von denen ich einige an dem Beispiel von $\infty^{2.3}$ aussprechen werde.

Die Mannigfaltigkeit von $2.\;3$ Dimensionen, die durch die homogene Gleichung
\[ f(x_1, x_2, x_3, x_4, x_5) = o \]
gegeben ist, werde durch die rationalen Substitutionsformeln
\[ (\varrho x_i = \varphi_i(y_1, y_2, y_3, y_4, y_5), \quad (i = 1, 2, \ldots 5) \]
in die Mannigfaltigkeit
\[ F(y_1, y_2, y_3, y_4, y_5) = o \]
übergeführt, und zwar so, dass sich auch umgekehrt mit Hülfe von $F = o$ die $y$ als rationale Functionen der $x$ ausdrücken lassen. Im Allgemeinen entsprechen sich dann $f = o$ und $F = o$ Punkt für Punkt eindeutig. Im Besondern kann es jedoch eintreten, dass einem Punkte der einen Mannigfaltigkeit eine Fläche $(\infty^{2.2})$ oder Curve $(\infty^{2.1})$ der andern, oder einer Curve der einen eine Fläche der andern entspricht.

\emph{Theorem 2.:} „Im Falle die Functionen $\varphi$ in einem einfachen Punkte von $F = o$

\origpage{302}
sämmtlich verschwinden, entspricht diesem Punkte auf $f = o$ eine Fläche, die sich rational durch zwei Parameter, oder eine Curve, die sich rational durch einen Parameter ausdrücken lässt.“

\emph{Theorem 3.:} „Im Falle die Functionen $\varphi$ längs einer einfachen Curve von $F = o$ sämmtlich verschwinden, entspricht dieser Curve auf $f = o$ eine Fläche, welche durch die Bewegung einer Curve erzeugt wird, die sich rational durch einen Parameter ausdrückt. Diese Fläche hat das Flächengeschlecht $p = o$, lässt sich aber im Allgemeinen \emph{nicht} rational durch zwei Parameter ausdrücken.

Die Flächen von höherm Geschlecht $p$ kann man immer vielfachen Curven auf $F = o$ entsprechen lassen.

\section*{III.}

Wenn eine Fläche die Eigenschaft hat, dass sich ihre Coordinaten in folgender Weise darstellen lassen:
\[ \varrho x_i = \varphi_i(\mu,\nu).\; \lambda^{r} + \varphi'_i(\mu\nu) \lambda^{r-1} + \ldots + \varphi_i^{(\sigma)}(\mu,\nu), \tag{1} \]
wobei
\[ \psi(\mu,\nu) = 0, \tag{2} \]
nämlich als rationale Functionen eines Parameters $\lambda$, deren Coefficienten rationale Functionen zweier Parameter $\mu, \nu$ sind, zwischen denen die von $\lambda$ unabhängige algebraische Gleichung (2) besteht, so hat die Fläche das Flächengeschlecht $p = o$.

Alle Flächen mit $p = o$, die wenigstens \emph{eine} einfach unendliche Schaar rationaler Curven enthalten, lassen sich in die Form (1), (2) setzen. Diese Flächen unterscheiden sich

\origpage{303}
aber noch nach der Irrationalität der Curvengleichung (2), auf welche sie führen.

\emph{Definition.} Ich bezeichne als das \emph{Curvengeschlecht} $p'$ einer Fläche mit $p = o$ das Geschlecht der Curve (2), auf welche diese Fläche führt.

\emph{Theorem 4.:} „Die nothwendige und hinreichende${}^{1)}$ Bedingung für das eindeutige Entsprechen zweier Flächen mit $p = o$ ist die Gleichheit des Curvengeschlechts $p'$ der beiden Flächen.“

Hiernach zerfällt die Flächenklasse $p = o$ in weitere Klassen nach dem Curvengeschlecht $p'$. Dieses $p'$ ist das Geschlecht der ebenen Schnittcurve des Kegels, der durch Gleichung (2) dargestellt werden kann und der der gegebenen Fläche eindeutig, Punkt für Punkt entspricht. Man kann sich ferner die Fläche erzeugt denken durch irrationale Fortschreitung einer rationalen Curve, wobei die irrationale Fortschreitung das Curvengeschlecht bestimmt. Nach Theorem 3. ist dann $p'$ zugleich das Geschlecht der auf $F = o$ liegenden Curve, welcher die Fläche auf $f = o$ entspricht. Die Normalflächen dieser Klasse sind die über den Normalcurven${}^{2)}$ beschriebenen Kegel. Die Coordinaten der Flächen drücken sich aus als rationale Functionen eines Parameters, deren Coefficienten Abelsche Functionen sind. Es lässt sich daher die von Herrn Clebsch gegebene Anwendung dieser Functionen auf die Geometrie der Curven${}^{3)}$ auch auf die

\textbf{1)} Das Wort hinreichend bezieht sich natürlich nur auf die Form der beiden Gleichungen, die sich noch durch die speciellen Werthe der Moduln von einander unterscheiden können.

\textbf{2)} Clebsch u.\ Gordan, Abel'sche Funct.\ §.\ 18.

\textbf{3)} Crelle's Journal, Bd.\ 63.

\origpage{304}
Geometrie der Flächen vom Flächengeschlecht $o$ übertragen.

Unter diese Form fallen vor Allem die \emph{windschiefen Flächen}, die alle das Flächengeschlecht $= o$ besitzen, während sie durch ihr Curvengeschlecht in Klassen zerfallen. Ihr Curvengeschlecht ist zugleich das ihres ebenen Querschnitts.${}^{1)}$

Als Corollar hat man noch den für die Abbildung von Flächen auf Ebenen wichtigen Satz: „das nothwendige und hinreichende Kriterium der Abbildbarkeit einer Fläche auf einer Ebene ist, dass die Fläche das Flächengeschlecht $p$ und das Curvengeschlecht $p' = o$ hat.“

Für Flächen von höherem Flächengeschlecht, sowie für höhere Mannigfaltigkeiten lassen sich ähnliche Kriterien und Klasseneintheilungen aufstellen, die man am Einfachsten aus den dem Theorem 3) analogen Theoremen ableitet. Ich werde die nähere Ausführung, sowie die Beweise der vorstehenden Theoreme an einem anderen Orte mittheilen.

\section*{IV.}

Man kann von dieser Theorie eine für die Geometrie wichtige Anwendung auf die Liniencomplexe ersten und zweiten Grades machen. Die Elimination einer der Coordinaten aus den beiden einem der Complexe entsprechenden Gleichungen ergiebt eine in 5 Coordinaten homogene Gleichung. Diese stellt ein höheres Gebilde dar, das man ansehen kann als erzeugt durch die rationale Fortbewegung einer Fläche, die man auf einer Ebene abbilden kann. Hieraus

\textbf{1)} Diese Eintheilung der Regelflächen stimmt mit der von Herrn Schwarz in Crelle's Journal Bd.\ 64 und 67, und von Herrn Lüroth Bd.\ 67 gegebenen überein.

\origpage{305}
folgt, dass jenes Gebilde sowohl das allgemeine Geschlecht $p$, als das Flächen- und Curvengeschlecht $= o$ hat. Daher ist es möglich, die Liniencomplexe des ersten und zweiten Grades so in unserm \emph{Raum} abzubilden, dass im Allgemeinen jedem Punkt des Raumes \emph{ein} Strahl des Complexes und umgekehrt entspricht, und man kann durch eindeutige Abbildung die ganze Geometrie unseres Raumes auf die Geometrie dieser Complexe übertragen.

Wenn die Gleichungen eines Liniencomplexes ersten Grads in der Form geschrieben werden:
\[
\begin{aligned}
o &= x_1 x_4 + x_2 x_5 + x_3 x_6 \\
o &= \Phi(x_1 x_2, x_3, x_4) + A x_5 + B x_6
\end{aligned}
\tag{3}
\]
wo $\Phi$ eine lineare Function, $A$ und $B$ Constanten sind, so wird die einfachste Abbildung durch die Formeln gegeben:
\[
\begin{aligned}
\varrho x_1 &= \xi_1 (A\xi_3 - B\xi_2) \\
\varrho x_2 &= \xi_2 (A\xi_3 - B\xi_2) \\
r x_3 &= \xi_3 (A\xi_3 - B\xi_2) \\
\varrho x_4 &= \xi_4 (A\xi_3 - B\xi_2) \\
\varrho x_5 &= \xi_1 \xi_4 B - \xi_3 \Phi(\xi_1, \xi_2, \xi_3, \xi_4) \\
\varrho x_6 &= - \xi_1, \xi_4 A + \xi_2 \Phi(\xi_1, \xi_2, \xi_3, \xi_4)
\end{aligned}
\tag{4}
\]
wobei das im Raume liegende Fundamentalgebilde aus einem Kegelschnitt besteht.

Die einfachste Abbildung eines Liniencomplexes zweiten Grads verdanke ich Herrn Dr.\ Klein. Wenn man den Complex auf ein Tetraeder bezieht, welches so gegen ihn liegt, dass alle geraden Linien, welche durch den den beiden Kanten 5 und 6 gemeinsamen Punkt innerhalb der durch dieselben bestimmten Ebene hindurchgehen, dem Complex angehören, so lassen sich die Gleichungen desselben ebenfalls in der Form (3) schreiben, wobei man nur unter $\Phi$ eine homogene Function zweiten Grades, unter $A$ und $B$ lineare homogene Functionen

\origpage{306}
von $x_1, x_2, x_3, x_4$ zu verstehen hat. Denkt man sich noch in den Formeln (4) in $A, B$ an Stelle von $x_1, x_2, x_3, x_4$ resp.\ $\xi_1, \xi_2, \xi_3, \xi_4$ gesetzt, so wird der Complex im Raume abgebildet durch diese Functionen dritter Ordnung, und als Fundamentalgebilde tritt eine Curve 5.\ Ordnung auf.

\emph{Göttingen}, den 30.\ Juni 1869.

\end{document}
